\section{Introduction}
\glsresetall

Large-scale experimental datasets have become increasingly important for wireless communication, sensing, and positioning research because they decouple algorithm development from scarce hardware testbed time and expose downstream users to realistic synchronization errors, coverage gaps, and multimodal alignment constraints. Recent multimodal releases such as DeepSense~6G show how carefully documented measurement assets can accelerate communication-aware sensing research well beyond their original experimental purpose~\cite{alkhateeb2023deepsense}. In parallel, Techtile was designed exactly as the kind of room-scale, distributed, synchronized infrastructure needed to create such assets for beyond-5G and 6G experimentation~\cite{techtilePrimer,techtileOpen6g}.

\dataset{} packages one concrete synchronized measurement campaign on top of Techtile. It includes the acquisition stack, rover control, acoustic tooling, \gls{rf} measurement scripts, parsing code, completeness checks, processed data products, and notebook-based analysis workflow used during the campaign. This makes the release more than a collection of files: it is the operational record of how synchronized rover, acoustic, and \gls{rf} measurements were acquired and transformed into analysis-ready products.

The closest dataset families motivate, but do not duplicate, this measurement design. DeepSense~6G provides a broad multi-scenario resource for communication-aware multimodal sensing, with RF measurements aligned to exteroceptive sensing modalities such as camera, radar, LiDAR, and position information~\cite{alkhateeb2023deepsense}. Its scale and scenario diversity are complementary to the present release. \dataset{} instead targets an indoor distributed positioning setting in which each acquisition cycle joins a ceiling-aperture \gls{rf} snapshot, an ultrasonic room response, and a motion-capture ground-truth pose at the same rover stop.

Indoor \gls{csi} localization work has established the utility of fingerprinting and learning-based position inference from channel measurements~\cite{kazemi2022csifingerprint,kordi2024survey,dai2023wifi}. Massive-MIMO and distributed-MIMO studies further show that large or spatially distributed RF apertures can provide richer spatial structure for fine-grained localization and data-driven channel modeling~\cite{li2022massivemimo,euchner2022bigcsidata}. These resources are essential for RF-only positioning research, but they generally do not provide synchronized acoustic waveforms recorded at the same physical pose. Conversely, acoustic positioning datasets and Techtile acoustic studies emphasize time-of-arrival, room impulse response, and microphone-array processing in the audible or ultrasonic domain~\cite{techtileAcoustic2022}. They typically do not include phase-coherent distributed RF channel observations that can be fused cycle by cycle with the acoustic traces.

This places \dataset{} between three established lines of work: multimodal sensing datasets, RF \gls{csi} localization datasets, and acoustic positioning datasets. Its distinguishing element is not merely that several files are released, but that the same orchestrator cycle defines the rover pose, acoustic capture, and distributed \gls{rf} measurement. The cycle-level key makes the RF and acoustic products directly joinable, while Techtile's shared synchronization and calibration infrastructure makes the distributed RF observations phase-coherent across the active ceiling aperture~\cite{syncWptTechtile}. The release also follows open-science dataset practice by exposing processing scripts, completeness checks, NetCDF products, and notebooks rather than only final figures or summary metrics~\cite{dss6g}.

Relative to earlier Techtile infrastructure and acoustic-only work, \dataset{} contributes a dense scanned campaign with a co-located RF-acoustic rover payload rather than only platform characterization or acoustic measurements. Relative to DeepSense~6G, the distinguishing modality pairing is synchronized \gls{rf}+acoustic propagation inside a controlled indoor distributed aperture rather than broad outdoor RF-perception sensing. Relative to pure \gls{csi} datasets, the distinguishing capability is cycle-level synchronization of phase-coherent distributed \gls{rf} measurements with raw acoustic waveforms, common ground truth, and an open processing workflow.

This paper is intentionally not framed as a benchmark paper for a specific estimator, classifier, or neural architecture. Its purpose is narrower and more practical: to document the experiments that were actually run, quantify what is currently complete in the released data products, and explain what \dataset{} already enables a researcher to do. In that sense, the paper serves as a rigorous companion to the data release. It reconstructs the measurement workflow, summarizes the campaign experiment by experiment, clarifies how the NetCDF products should be interpreted, and distinguishes between the total logged campaign and the presently usable tri-modal subset.

Three aspects of \dataset{} are especially worth documenting. First, the acquisition is genuinely synchronized at the level of the physical rover stop: one orchestrator cycle corresponds to one commanded rover position, one Qualisys sample, one acoustic capture, and one \gls{rf} cycle. Second, the processed products are already structured for downstream use, with explicit availability masks and notebook-backed analysis examples. Third, completeness is uneven across experiments, which makes it essential to state clearly what parts of the campaign support \gls{rf}-only analysis and what parts support fully aligned \gls{rf}+acoustic+position analysis.

The remainder of the paper is as follows. \Cref{sec:platform} describes the platform and acquisition workflow. \Cref{sec:campaign} quantifies the campaign and its present coverage. \Cref{sec:capabilities} discusses the recorded data products and the positioning capabilities already exposed in the analysis workflow. \Cref{sec:discussion} concludes with the main interpretation caveats and current limitations.
